\section*{\texorpdfstring{\S\CurrentExerciseGroup}{第{\CurrentExerciseGroup}节}}
\phantomsection
\addcontentsline{toc}{section}{第{\CurrentExerciseGroup}节习题}

\begin{enumerate}
\def\labelenumi{\arabic{enumi})}
\tightlist
\item
  设 \(\mathbf{T}\) 与 \(\mathbf{X}\) 为两个局部紧空间，\(\mu\) 为 \(\mathbf{T}\) 上的正测度，\(\pi\) 为一个连续且 \(\mu\) -逆紧的映射，把 \(\mathbf{T}\) 映入 \(\mathbf{X}\) 中，并令 \(\nu = \pi(\mu)\)。
\end{enumerate}

\begin{enumerate}
\def\labelenumi{\alph{enumi})}
\item
  证明：对每个数值函数 \(g \geqslant 0\)（在 \(X\) 上定义且下半连续），\(g \circ \pi\) 在 \(T\) 上下半连续，且 \(\nu^{*}(g) = \mu^{*}(g \circ \pi)\)。
\item
  证明：若 \(\mathbf{f}\) 是取值于一个 Banach 空间的 \(\nu\) -可积函数，则 \(\mathbf{f} \circ \pi\) 是 \(\mu\) -可积的，且 \(\int \mathbf{f} d\nu = \int (\mathbf{f} \circ \pi) d\mu\)。在没有附加条件时逆命题是否成立（参见 §4, 习题 2b)）？
\end{enumerate}

¶ 2) 设 T 与 X 为两个局部紧空间，P 为 T 到 X 中的逆紧连续映射所成之集，赋以紧收敛拓扑。设 H 为 P 的一个子集，使对 X 的每个紧子集 K，诸集 \(\overline{\pi}^{-1}(K)\)（其中 \(\pi\) 跑遍 H）的并在 T 中相对紧。

\begin{enumerate}
\def\labelenumi{\alph{enumi})}
\item
  证明：映射 \((\pi, \lambda) \mapsto \pi(\lambda)\)（从 \(\mathrm{H} \times \mathcal{M}(\mathrm{T})\) 到 \(\mathcal{M}(\mathrm{X})\) 中）当空间 \(\mathcal{M}(\mathrm{T})\) 与 \(\mathcal{M}(\mathrm{X})\) 都赋以拟强拓扑时未必连续（Ch. III, §1, 习题 8）。
\item
  若 B 为 \(\mathcal{M}(\mathrm{T})\) 的一个有界子集，证明 \(\mathrm{H} \times \mathrm{B}\) 上的限制（即映射 \((\pi, \lambda) \mapsto \pi(\lambda)\) 的限制）当空间 \(\mathcal{M}(\mathrm{T}), \mathcal{M}(\mathrm{X})\) 中每一个都赋以淡拓扑时是连续的（利用 Ch. III, §1 的命题 15）。
\item
  证明：映射 \((\pi, \lambda) \mapsto \pi(\lambda)\)（从 \(\mathrm{H} \times \mathcal{M}_{+}(\mathrm{T})\) 到 \(\mathcal{M}_{+}(\mathrm{X})\) 中）当 \(\mathcal{M}(\mathrm{T})\) 与 \(\mathcal{M}(\mathrm{X})\) 赋以淡拓扑时连续（参见 Ch. III, §1, 习题 10）。 d) 给出一个例子，其中 \(\mathrm{T} = \mathrm{X}\) 为紧的，且映射 \((\pi, \lambda) \mapsto \pi(\lambda)\)（从 \(\mathcal{P} \times \mathcal{M}(\mathrm{T})\) 到 \(\mathcal{M}(\mathrm{X}) = \mathcal{M}(\mathrm{T})\) 中）当 \(\mathcal{M}(\mathrm{T})\) 赋以淡拓扑时不连续（取 \(\mathrm{T}\) 为环面 \(\mathbf{T}\)，并利用 Ch. III, §4 的习题 2 b)）。
\end{enumerate}

¶ 3) 设 T 为非紧的局部紧空间，\(\mu\) 为 T 上的正测度，\(\mathcal{G}\) 为 T 的同胚所成的群。给出一个例子说明：映射 \(\pi \mapsto \pi(\mu)\)（从 \(\mathcal{G}\) 到 \(\mathcal{M}(T)\) 中）当 \(\mathcal{G}\) 赋以紧收敛拓扑且 \(\mathcal{M}(T)\) 赋以淡拓扑时未必连续（取 T 为 R 中由点 0 与点 n 和 1/n（n 为整数 \(\geqslant 1\)）组成的子空间，取 \(\mu\) 为由 \(\mu(\{0\}) = 0\)，\(\mu(\{1/n\}) = 1/n^2\)，\(\mu(\{n\}) = 1\) 定义的测度）。

\begin{enumerate}
\def\labelenumi{\arabic{enumi})}
\setcounter{enumi}{3}
\item
  对每个局部紧空间 T，设 \(\mathcal{M}^{c}(T)\) 为 \(\mathcal{M}(T)\) 中由具有紧支集的实测度组成的子空间。若 T 与 X 为两个局部紧空间，且 \(\pi\) 是 T 到 X 中的连续映射，则 \(\pi\) 是 \(\lambda\) -逆紧的，且 \(\pi(\lambda) \in \mathcal{M}^{c}(X)\) 对每个测度 \(\lambda \in \mathcal{M}^{c}(T)\) 成立。给出一个例子说明：映射 \(\lambda \mapsto \pi(\lambda)\)（从 \(\mathcal{M}^{c}(T)\) 到 \(\mathcal{M}^{c}(X)\) 中）当 \(\mathcal{M}^{c}(T)\) 与 \(\mathcal{M}^{c}(X)\) 赋以淡拓扑（或拟强拓扑（Ch. III, §1, 习题 8））时未必连续，且 \(\pi\) 不是逆紧映射（参见 Ch. III, §2, 习题 1）。
\item
  \begin{enumerate}
  \def\labelenumii{\alph{enumii})}
  \tightlist
  \item
    设 \(\rho\) 与 \(\sigma\) 为 \(\mathbf{R}\) 上的两个正测度。证明：若 \(\rho(]a,b]) = \sigma(]a,b])\) 对每个半开区间 \(]a,b]\) 成立，则 \(\rho = \sigma\)。
  \end{enumerate}
\end{enumerate}

\begin{enumerate}
\def\labelenumi{\alph{enumi})}
\setcounter{enumi}{1}
\tightlist
\item
  设 I 为 \(\mathbf{R}\) 的一个区间，左端点为 \(\alpha\)，右端点为 \(\beta\)，并设 \(\psi\) 为定义在 I 上的递增有限数值函数；以任意方式把 \(\psi\) 延拓到 \(\mathbf{R}\)。为使 \(\psi\) 对测度 \(\varphi_{\mathrm{I}} \cdot \mu\)（\(\mu\) 为 \(\mathbf{R}\) 上的 Lebesgue 测度）是逆紧的，必须且只需下列两个条件成立：\(1^{\circ}\) 或者 \(\alpha\) 有限，或者 \(\alpha = -\infty\) 且 \(\lim_{x \to -\infty} \psi(x) = -\infty\)；\(2^{\circ}\) 或者 \(\beta\) 有限，或者 \(\beta = +\infty\) 且
\end{enumerate}

\(\lim_{x\to +\infty}\psi (x) = +\infty .\)

假设这些条件已验证。对每个 \(y \in \mathbf{R}\)，设 \(\theta(y)\) 为 \(x \in I\) 中使 \(\psi(x) \leqslant y\) 成立者所成之集的上确界（若该集为空，则上确界等于 \(\alpha\)）。证明 \(\theta\) 是 \(\mathbf{R}\) 中的有限数值函数，递增且右连续。若 \(\nu\) 为 \(\varphi_{\mathrm{I}} \cdot \mu\) 在 \(\psi\) 下的像，则 \(\nu(]a,b]) = \theta(b) - \theta(a)\) 对每个半开区间 \(]a,b]\)（在 \(\mathbf{R}\) 中）成立。

\begin{enumerate}
\def\labelenumi{\alph{enumi})}
\setcounter{enumi}{2}
\tightlist
\item
  反之，证明对每个有限数值函数 \(\theta\)（在 \(\mathbf{R}\) 中递增且右连续），存在唯一的一个正测度 \(\nu\)（在 \(\mathbf{R}\) 上），使 \(\nu(]a,b]) = \theta(b) - \theta(a)\) 对每个半开区间 \(]a,b]\)（属于 \(\mathbf{R}\)）成立（“\(\mathbf{R}\) 上由 \(\theta'\) 定义的 Stieltjes 测度”；用 \(\int f d\theta\) 代替 \(\int f d\nu\) 来记）；此外，\(\mathbf{R}\) 上的每个正测度都可这样得到。在什么条件下测度 \(\nu\) 是弥散的？此时 \(\nu\) 在 \(\theta\) 下的像是什么？
\end{enumerate}

\begin{enumerate}
\def\labelenumi{\arabic{enumi})}
\setcounter{enumi}{5}
\tightlist
\item
  \begin{enumerate}
  \def\labelenumii{\alph{enumii})}
  \tightlist
  \item
    设 K 为区间 {}0,1{} 中的 Cantor 三分集（GT, IV, §2, No.~5）。证明存在一个弥散正测度 \(\nu\)（在 \(\mathbf{R}\) 上，以 K 为支集且总质量为 1），以及一个连续递增函数 \(\theta\)（定义于 \(\mathbf{R}\)），使 \(\theta(\mathbf{R}) = \mathbf{I}\) 且 \(\theta(\nu) = \varphi_{\mathrm{I}} \cdot \mu\)（\(\mu\) 为 Lebesgue 测度；参见 GT, IV, §8, 习题 16）。
  \end{enumerate}
\end{enumerate}

\begin{enumerate}
\def\labelenumi{\alph{enumi})}
\setcounter{enumi}{1}
\tightlist
\item
  由 a) 推出：存在 R 上的一个弥散正测度，它与 Lebesgue 测度互奇异，且其支集为整个 I（在与 K 相邻且含于 I 的每个区间 J 中，取一个与 \(\nu\) 在 I 到 J 上的仿射映射下的像成比例的测度；然后递归地进行）。
\end{enumerate}

\begin{enumerate}
\def\labelenumi{\arabic{enumi})}
\setcounter{enumi}{6}
\tightlist
\item
  设 \(\mu\) 为 \(\mathbf{R}\) 上的 Lebesgue 测度，I 为区间 \([0,1]\)（\(\mathbf{R}\) 中的）。若令 \(\theta(x) = |x|\)，则 \(\theta\) 对测度 \(\varphi_{\mathrm{I}} \cdot \mu\) 是逆紧的，且 \(\theta(\varphi_{\mathrm{I}} \cdot \mu) = \varphi_{\mathrm{I}} \cdot \mu\)。给出一个对 \(\varphi_{\mathrm{I}} \cdot \mu\) 可忽略的集的例子，使它在 \(\theta\) 下的像对 \(\varphi_{\mathrm{I}} \cdot \mu\) 不可测（参见 Ch. IV, §4, 习题 8）。
\end{enumerate}

¶ 8) 设 T 为紧空间，μ 为 T 上总质量为 1 的弥散正测度。

\begin{enumerate}
\def\labelenumi{\alph{enumi})}
\tightlist
\item
  证明存在一个连续映射 \(\pi\)（从 T 到 E = {}0,1{} 上），使 \(\mu\) 在 \(\pi\) 下的像是 E 上的 Lebesgue 测度 \(\lambda\)。（按 Urysohn 定理（GT, IX, §4, 定理 1）证明的方式进行：对每个 \(t\)（其中 \(0 \leqslant t \leqslant 1\)），定义一个开集 \(\mathrm{U}(t) \subset \mathrm{T}\)（对 \(\mu\) 可求积（Ch. IV, §5, 习题 17）），使 \(\mathrm{U}(0) = \varnothing\)，\(\mathrm{U}(1) = \mathrm{T}\)，\(\overline{\mathrm{U}(t)} \subset \mathrm{U}(t')\)（当 \(t < t'\) 时），最后 \(\mu(\mathrm{U}(t)) = t\)；利用 §5 的习题 7 证明：若 V、W 为 T 中两个可求积开集，使 \(\overline{\mathrm{V}} \subset \mathrm{W}\) 且 \(\mu(\mathrm{V}) < \mu(\mathrm{W})\)，则存在一个可求积开集 U 使 \(\overline{\mathrm{V}} \subset \mathrm{U} \subset \overline{\mathrm{U}} \subset \mathrm{W}\) 且使
\end{enumerate}

\[
\frac {1}{3} \mu (\mathrm{W-V}) \leqslant \mu (\mathrm{U-V}) \leqslant \frac {2}{3} \mu (\mathrm{W-V}).
\]

最后，应用习题 5 a)）。

\begin{enumerate}
\def\labelenumi{\alph{enumi})}
\setcounter{enumi}{1}
\item
  由 a) 推出：存在 T 的不是 \(\mu\) -可测的子集（Ch. IV, §4, 习题 8）；存在 \(\mathcal{L}^{p}(\mathrm{T},\mu)\) 中的序列，它 p 阶平均收敛于 0 但在 T 的任何点都不收敛于 0（\(1 \leqslant p < +\infty\)；Ch. IV, §3, 习题 1）；还存在序列 \((f_{n})\)，使 \((\widetilde{f}_{n})\) 对 \(\mathrm{L}^{p}(\mathrm{T},\mu)\) 的弱化拓扑收敛于 0 但不依测度收敛于 0（§5, 习题 19）。
\item
  再假设 T 可度量化。证明存在 T 的 \(\mu\) -可忽略子集 N、E 的 \(\lambda\) -可忽略子集 M，以及 E−M 到 T−N 上的一个同胚 \(\pi\)，使得：若把 \(\pi\)（任意地）延拓为 E 到 T 中的一个映射 \(\varphi\)，并把 \(\pi^{-1}\) 延拓为 T 到 E 中的一个映射 \(\psi\)，则 \(\varphi(\lambda)=\mu\) 且 \(\psi(\mu)=\lambda\)。（利用 Ch. IV, §5 的习题 17 证明：对每个整数 n \textgreater 0，存在 T 的一个有限分割，由一个 \(\mu\) -可忽略集以及一些可求积开集组成，其直径 \(\leqslant1/n\)（对与 T 的拓扑相容的度量而言）且测度 \(\leqslant1/n\)。递归地进行，然后通过取极限，由此推出存在 E−D 到 T 中的一个连续映射 f，其中 D 是 E 的可数子集，使 \(f(\lambda)=\mu\)；证明可以做到使 f 是 E−D 到 T 中一个测度为 1 的子集上的同胚（参见 §8, 习题 14）。）
\end{enumerate}

¶ 9) 设 \(\mu\) 与 \(\nu\) 为局部紧空间 T 上的两个弥散正测度。证明：为使 \(\nu\) 是以 \(\mu\) 为基的测度，只需 T 的每个 \(\mu\) -可测子集也是 \(\nu\) -可测的（用反证法，利用 §5 的定理 3 与命题 13，以及 §6 的习题 8 b)）。

\begin{enumerate}
\def\labelenumi{\arabic{enumi})}
\setcounter{enumi}{9}
\tightlist
\item
  在区间 \(\mathbf{I} = ]0,1[\)（含于 \(\mathbf{R}\)）中定义函数 \(g\) 为 \(g(t) = -\sqrt{n}\)（当 \(\frac{1}{2}\left(\frac{1}{n} + \frac{1}{n+1}\right) < t \leqslant \frac{1}{n}\) 时）与 \(g(t) = \sqrt{n}\)（当 \(\frac{1}{n+1} < t \leqslant \frac{1}{2}\left(\frac{1}{n} + \frac{1}{n+1}\right) (n = 1,2,\ldots)\) 时），并把 \(g\) 延拓为在 \(\mathbf{R} - \mathbf{I}\) 上取 0；\(g\) 对 Lebesgue 测度可积。令
\end{enumerate}

\[
\mathrm{G} (x) = \int_ {0} ^ {x} g (t) d t;
\]

证明：函数 \(f\)（等于 \(1 / \sqrt{x}\)（在 \([-1, +1]\) 上），在其外等于 0，且对 Lebesgue 测度可积）使得 \(t \mapsto f(\mathrm{G}(t))g(t)\) 在 I 上不可积。

¶ 11) 记号沿用 No.~5 的；g 是 I 上的一个局部 μ-可积数值函数。

\begin{enumerate}
\def\labelenumi{\alph{enumi})}
\item
  为使 G 是 I 到 G(I) 中的 \(\lambda\) -逆紧映射，必须且只需下列条件成立：\(1^{\circ}\) 极限 G(a+) 与 G(b−) 在 \(\overline{\mathbf{R}}\) 中存在；\(2^{\circ}\) 若 G(a+) ∈ G(I)，则 g 是 \(\mu\) -可积的（在区间 {]} \(a, x_0[\) 上）；\(3^{\circ}\) 若 G(b−) ∈ G(I)，则 g 是 \(\mu\) -可积的（在区间 {]} \(x_0, b[\) 上）。
\item
  假设极限 \(\mathrm{G}(a+)\) 与 \(\mathrm{G}(b-)\) 在 \(\overline{R}\) 中存在。证明：若 f 使 \(t \mapsto \mathbf{f}(\mathrm{G}(t))g(t)\) 在 I 上对 Lebesgue 测度可积，则 f 在以 \(\mathrm{G}(a+)\) 与 \(\mathrm{G}(b-)\) 为端点的区间上对 Lebesgue 测度可积，且公式 (13) 成立。（化归于 \(x_{0}\) 是数 a、b 之一的情形；例如若 \(x_{0} = a\)，则注意 I 中存在一个严格递增且趋于 b 的序列 \((b_{n})\)，使序列 \((\mathrm{G}(b_{n}))\) 递增或递减；应用 Ch. IV, §4, No.~3 的命题 4 与 Lebesgue 定理。）
\end{enumerate}

¶ 12) a) 设 g 为定义在 R 的紧区间 {[}a, b{]} 上的一个有限数值函数。称 g 在 I 中的右展集（相应地，左展集）为这样的 \(x \in I\) 所成之集：存在 \(y \in I\)，使 x < y（相应地，x \textgreater{} y）且 \(g(x) < g(y)\)。证明：若 g 连续，则 g 在 I 中的右（相应地，左）展集在 I 中是开的，并且，若 {]} \(\alpha, \beta[\) 是该集的一个连通分支 \(^{1}\)，则 \(g(\alpha) = g(\beta)\)，且 \(g(x) \leqslant g(\alpha)\)（对 \(\alpha < x < \beta\)）。

\begin{enumerate}
\def\labelenumi{\alph{enumi})}
\setcounter{enumi}{1}
\item
  设 \(r_1, r_2\) 为使 \(0 \leqslant r_1 < r_2\) 的两个实数。设 \(g\) 为 I 上的一个递增连续函数；设 \(E'\) 为 \(g(x) - r_1x\) 在 I 中的左展集，设 \(E''\) 为 \(g(x) - r_2x\) 在 \(E'\) 的每个连通分支的闭包所成各个区间中的右展集的并。若 \(\mu\) 为 I 上的 Lebesgue 测度，试利用 a) 证明 \(\mu(E'') \leqslant \frac{r_1}{r_2} \mu(E')\)。
\item
  称在点 \(x \in I\) 处，有限数值函数 \(f\)（定义在 \(I\) 上）的右上导数与右下导数为下列各数
\end{enumerate}

\[
\begin{array}{l} \mathrm{D} _ {r} ^ {+} f (x) = \operatorname * {l i m s u p} _ {h \to 0,   h > 0} \left(f (x + h) - f (x)\right) / h \\ \mathrm{D} _ {r} ^ {-} f (x) = \operatorname * {l i m i n f} _ {h \to 0,   h > 0} \left(f (x + h) - f (x)\right) / h. \end{array}
\]

类似地，称 \(f\) 在点 \(x\) 处的左上导数与左下导数为下列各数

\[
\begin{array}{l}\mathrm{D} _ {l} ^ {+} f (x) = \underset {h \rightarrow 0, h <   0} {\lim \sup} \left(f (x + h) - f (x)\right) / h\\\mathrm{D} _ {l} ^ {-} f (x) = \underset {h \rightarrow 0, h <   0} {\lim \inf} \left(f (x + h) - f (x)\right) / h.\end{array}
\]

  证明：若 \(f\) 连续且递增，则使 \(x \in I\) 且 \(\mathrm{D}_r^+ f(x) = +\infty\) 的点所成之集与使 \(x \in I\) 且 \(\mathrm{D}_l^- f(x) < \mathrm{D}_r^+ f(x)\) 的点所成之集对 \(\mu\) 的测度为零（对每个有理数 \(r > 0\)（相应地，对每对有理数 \((r_1, r_2)\)，其中 \(0 \leqslant r_1 < r_2\)），考察使 \(x \in I\) 且 \(\mathrm{D}_r^+ f(x) > r\) 的点所成之集（相应地，使 \(\mathrm{D}_l^- f(x) < r_1\) 与 \(\mathrm{D}_r^+ f(x) > r_2\) 同时成立的点），并应用 \(b\)））。由此推出，对几乎每个 \(x \in I\)，\(f\) 都有一个有限导数（“Lebesgue 定理”）（把前面的结果也应用于 \(-f(-x)\)）。

¶ 13) 在 R 的紧区间 {[}a, b{]} 中，设 \((f_n)\) 为递增的有限连续函数序列，使对一切 n 有 \(f_n(a) = 0\)，且和 \(s(x) = \sum_{n=1}^{\infty} f_n(x)\) 在 I 上有限。证明：存在一个集 N ⊂ I，对 Lebesgue 测度可忽略，使对每个 \(x \in I - N\)，导数 \(f_n'(x)\) 与 \(s'(x)\) 存在，且 \(s'(x) = \sum_{n=1}^{\infty} f_n'(x)\)（“Fubini 定理”）。（利用习题 12。注意以 \(f_n'(x)\) 为通项的级数几乎处处收敛；令

\[
s _ {n} (x) = \sum_ {k = 1} ^ {n} f _ {k} (x),
\]

接着考察一个子序列 \((s_{n_k})\)，使以 \(s(x) - s_{n_k}(x)\) 为通项的级数在 I 上收敛，并把前面的注记应用于该级数。）

\begin{enumerate}
\def\labelenumi{\arabic{enumi})}
\setcounter{enumi}{13}
\item
  设 \(f\) 为定义在紧区间 \(I = [a, b]\)（\(\mathbf{R}\) 中）上的数值函数，在 \(I\) 上对 Lebesgue 测度可积。若令 \(F(x) = \int_{a}^{x} f(t) dt\)，证明 \(F\) 在 \(I\) 中几乎处处有一个等于 \(f\) 的导数（化归于 \(f \geqslant 0\) 的情形；然后，利用习题 13，相继考察 \(f\) 下半连续的情形与一般情形（参见 Ch. IV, §4, No.~4, 定理 3 的推论））。
\item
  用 \(\mu\) 表示 \(\mathbf{R}\) 上的 Lebesgue 测度。称 \(x \in \mathbf{R}\) 为 \(\mathbf{R}\) 的子集 A 的密度点，如果当 \(h\) 与 \(k\) 经取 \(>0\) 的值而趋于 0 时，商 \(\mu^{*}(\mathrm{A} \cap [x - h, x + k]) / (h + k)\) 趋于 1。证明：\(\mathbf{R}\) 的任意子集 A 中不是 A 的密度点的点所成之集对 \(\mu\) 可忽略。（化归于 A 含于紧区间 \([a, b]\) 的情形，并考察函数 \(s_{\mathrm{A}}(x) = \mu^{*}(\mathrm{A} \cap [a, x])\)。取一个包含 A 的开集递减序列 \((\mathrm{A}_n)\)，使以 \(s_{\mathrm{A}_n}(x) - s_{\mathrm{A}}(x)\) 为通项的级数在 \([a, b]\) 上收敛，并利用习题 13。）
\item
  \begin{enumerate}
  \def\labelenumii{\alph{enumii})}
  \tightlist
  \item
    设 \(g\) 为对 Lebesgue 测度可积的函数（在 \(\mathbf{E} = [0,1]\) 上）。令 \(\mathrm{G}(x) = \int_0^x g(t)dt\)；证明：在 \(x \in \mathbf{E}\) 中 \(g\) 连续的每一点处，\(\mathbf{G}\) 有一个等于 \(g\) 的导数。
  \end{enumerate}
\end{enumerate}

\begin{enumerate}
\def\labelenumi{\alph{enumi})}
\setcounter{enumi}{1}
\item
  给出函数 \(g\) 的一个例子：它有界且在 \(\mathbf{E}\) 中几乎处处连续，而 \(\mathbf{G}\) 在 \(\mathbf{E}\) 的一个不可数子集的各点处没有右导数（参见 FRV, I, §2, 习题 8）。
\item
  证明：等于 \(x^2 \sin \frac{1}{x^2}\)（当 \(x \neq 0\) 时）而当 \(x = 0\) 时等于 0 的函数在 \(E\) 的每一点处有导数，但这导数在 \(E\) 上对 Lebesgue 测度不可积。
\end{enumerate}

\begin{enumerate}
\def\labelenumi{\arabic{enumi})}
\setcounter{enumi}{16}
\item
  设 T 与 X 为两个局部紧空间，\(\pi\) 为 T 到 X 中的连续映射。证明：若 \(\pi\) 对 T 上的每个正测度 \(\mu\) 都是 \(\mu\) -逆紧的，则 \(\pi\) 是逆紧的。（用反证法，利用 GT, I, §10 的命题 7。）
\item
  在命题 4 b) 中，若省去 \(\pi'\) 连续这一假设，结论可能不成立。（取 \(T = R\)，\(T' = \overline{R}\)，\(T'' = R\)；取 \(\mu\) 为 Lebesgue 测度，\(\pi\) 为典则单射；令 \(\pi'(x) = x\)（对 \(x \in R\)），并令 \(\pi'(\pm \infty) = 0\)。）
\item
  设 \(\mathrm{T}, \mathrm{X}, \mathrm{Y}\) 为局部紧空间，\(\pi : \mathrm{T} \to \mathrm{X}\) 与 \(\pi' : \mathrm{X} \to \mathrm{Y}\) 为映射，\(\pi'' = \pi' \circ \pi\)，\(\mu\) 为 \(\mathrm{T}\) 上的实测度。可能发生：\(\pi\) 是 \(\mu\) -逆紧的，且 \(\pi'\) 是 \(\pi(\mu)\) -逆紧的，而不必 \(\pi''\) 是 \(\mu\) -逆紧的。（取 \(\mathrm{T} = \mathbf{R} \times \{0, 1\}\)，\(\mathrm{X} = \mathbf{R}\)，\(\mathrm{Y} = \{0\}\)，\(\pi = \mathrm{pr}_1\)；取 \(\mu\) 为在 \(\mathbf{R} \times \{0\}\)（相应地，\(\mathbf{R} \times \{1\}\)）——两者典则等同于 \(\mathbf{R}\)——上诱导 Lebesgue 测度（相应地，Lebesgue 测度的相反测度）的测度。）
\item
  设 \(\mathrm{T},\mathrm{X},\mathrm{Y}\) 为三个局部紧空间，\(\mu\) 为 \(\mathrm{T}\) 上的正测度，\(t\mapsto \lambda_t\) 为 \(\mu\) -适宜的映射，把 \(\mathrm{T}\) 映入 \(\mathcal{M}_{+}(\mathrm{X})\) 中，令 \(\nu = \int \lambda_t d\mu (t)\)，并设 \(\pi\) 为 \(\nu\) -逆紧映射（把 \(\mathrm{X}\) 映入 \(\mathrm{Y}\) 中）。假设下列假设之一成立：a) \(\mathrm{Y}\) 无穷远处可数且 \(\nu\) 是适度的；b) \(\mathrm{Y}\) 无穷远处可数且 \(\lambda_t\) 在 \(\mathrm{T}\) 中局部几乎处处有界。则 \(\pi\) 是 \(\lambda_t\) -逆紧的（在 \(\mathrm{T}\) 中局部几乎处处），映射 \(t\mapsto \pi (\lambda_t)\) 是 \(\mu\) -适宜的，且其积分等于 \(\pi (\mu)\)。
\item
  设 \(\mathrm{T},\mathrm{X},\mathrm{Y}\) 为三个局部紧空间，\(\pi : \mathrm{T} \to \mathrm{X}\)、\(\pi': \mathrm{X} \to \mathrm{Y}\) 为普遍可测映射。则 \(\pi' \circ \pi\) 普遍可测。（设 \(\mu\) 为 \(\mathrm{T}\) 上的正测度。为证 \(\pi' \circ \pi\) 是 \(\mu\) -可测的，化归于 \(\mathrm{T}\) 紧且 \(\pi\) 连续的情形，并利用 \(\pi'\) 对 \(\pi(\mu)\) 可测这一事实。）
\item
  设 \(X, Y\) 为两个紧空间，\(U\) 为 Banach 空间 \(\mathcal{C}(X; \mathbf{R})\) 到 Banach 空间 \(\mathcal{C}(Y; \mathbf{R})\) 中的连续线性映射，\(\pi : X \to Y\) 为连续映射。假设对每个 \(y \in Y\) 及每个函数 \(f \in \mathcal{C}(X; \mathbf{R})\)，
\end{enumerate}

\[
\inf _ {\pi (x) = y} f (x) \leqslant (U \cdot f) (y) \leqslant \sup _ {\pi (x) = y} f (x).
\]

证明：在这些条件下，存在一个淡连续映射 \(\sigma: y \mapsto \sigma_y\)（从 Y 到 \(\mathcal{M}_+(X)\) 中），使对每个 \(y \in Y\)，\(\sigma_y\) 由 \(\overline{\pi}^{-1}(y)\) 承载，且 \((U \cdot f)(\pi(x)) = \langle f, \sigma_{\pi(x)} \rangle\)（对一切 \(x \in X\)）成立；反之亦然。（考察转置 \(^t U\) 。）

\setcounter{exercisegroup}{7}
\edef\CurrentExerciseGroup{\arabic{exercisegroup}}
